\chapter{Stability Under Stochastic Decoding (Day 7)}
\label{chap:day07}

\section{Motivation}

Stability is one of the six metrics defined for evaluating black-box explainable-AI frameworks
\citep{arreche2024exai}: rerun the same input through the model multiple times and ask whether the
answer, and the region the answer is grounded in, stays put. Days 3--6 all ran Florence-2 exactly
once per sample. Stability is the first day to ask what happens when the same sample is run more
than once and the two runs are compared against each other --- and that question carries a trap the
plan flagged explicitly before any code was written against it.

\textbf{The vacuous-result warning.} Florence-2's default decoding configuration, used throughout
Chapters 3--6, is beam search with \texttt{num\_beams=3} and no sampling: deterministic decoding,
in the sense that the same image and prompt, run through the same model weights, produce
bit-identical output every time. Rerunning a deterministic pipeline three times and comparing the
three outputs would report ``100\% stable'' by construction, for every sample, regardless of
whether Florence-2's underlying judgment is actually robust to anything. That result would be
vacuous --- true of any deterministic function applied to a fixed input, and informative about
nothing. Taking this warning seriously up front, before treating any stability number as
meaningful, is what shapes this chapter's implementation choice and its first reported check.

\section{Implementation}

\texttt{metrics/stability.py}'s \texttt{run\_n\_times()} avoids the vacuous case by switching
decoding mode for the rerun comparison specifically: each sample is rerun 3$\times$ with
\texttt{do\_sample=True}, \texttt{num\_beams=1}, \texttt{temperature=0.7} --- genuine stochastic
sampling, not the deterministic beam search used for the answers reported elsewhere in this pilot.
The three sampled reruns are then compared pairwise to compute \texttt{answer\_agreement\_rate}
(do the reruns agree with each other on the final yes/no safety-rule answer) and two region-overlap
scores, both intersection-over-union (IoU) measures computed across the three reruns'
top-ranked-region boxes. The deterministic beam-search answer that Chapters 3--5 already logged in
\texttt{baseline\_predictions.csv} is brought back in only as a separate comparison point ---
\texttt{deterministic\_matches\_majority} --- never folded into the stability score itself,
preserving the distinction between ``how stable is the model under genuine resampling'' and ``does
the model's one official answer happen to agree with the resampled majority.''

\subsection{Confirming the result is not vacuous}

Before trusting any aggregate stability number, the sampled-rerun setup was checked directly
against the failure mode it was designed to avoid: if every sample showed perfect agreement across
all three sampled reruns, that would be indistinguishable from the vacuous deterministic case, and
would mean \texttt{temperature=0.7} sampling was not actually exercising any real variability in
Florence-2's output. It is not vacuous. Mean \texttt{answer\_agreement\_rate} across all 163
samples is \textbf{0.775}, not 1.0 --- roughly a quarter of rerun pairs genuinely disagree on the
final answer. That gap from 1.0 is the evidence that stochastic decoding is doing real work here,
and it is what licenses treating the rest of this chapter's numbers as a measurement of something
real rather than an artifact of comparing a deterministic function against itself.

\section{Results}

\texttt{scripts/07\_stability\_test.py} ran the rerun pipeline over all 163 pilot samples: three
sampled reruns each, plus the separate comparison against the existing deterministic answer.

\begin{center}
\begin{tabular}{lr}
\toprule
\textbf{Metric} & \textbf{Mean (163 samples)} \\
\midrule
\texttt{answer\_agreement\_rate} (pairwise answer match across 3 reruns) & 0.775 \\
\texttt{deterministic\_matches\_majority} (beam search vs.\ sampled majority) & 83.4\% \\
\texttt{top\_region\_overlap\_score} (IoU of area-ranked top-1 region) & \textbf{0.667} \\
\texttt{object\_region\_overlap\_score} (IoU of the rule's safety-object box) & \textbf{0.602} \\
\bottomrule
\end{tabular}
\end{center}

\texttt{deterministic\_matches\_majority} at 83.4\% says that Florence-2's one official,
deterministic answer agrees with the majority answer across the three stochastic reruns for most,
but not all, samples --- a separate and reassuring signal that the deterministic pipeline used
everywhere else in this pilot is not systematically an outlier relative to what the model would say
under resampling, even though stability itself, 0.775, shows real rerun-to-rerun disagreement on
roughly a quarter of samples.

\subsection{By assigned rule}

\begin{center}
\begin{tabular}{lrr}
\toprule
\textbf{Rule} & \textbf{answer\_agreement\_rate} & \textbf{object\_region\_overlap\_score} \\
\midrule
rule\_1 (hard hat) & 0.735 & 0.499 \\
rule\_2 (harness) & 0.744 & 0.533 \\
rule\_3 (guardrail) & 0.782 & 0.615 \\
rule\_4 (excavator proximity) & 0.893 & \textbf{0.922} \\
\bottomrule
\end{tabular}
\end{center}

The by-rule breakdown already shows a pattern worth flagging on first read: rule\_4's excavator
answers and object boxes are both noticeably more stable than rules 1--3's, and the gap is far
larger for \texttt{object\_region\_overlap\_score} (0.499 at the low end to 0.922 at the high end)
than for \texttt{answer\_agreement\_rate} (0.735 to 0.893). The remainder of this chapter is about
that wider gap, and about the fact that the headline \texttt{top\_region\_overlap\_score} ---
0.667, reported as a single aggregate above --- does not show it at all.

\section{Finding: A Hidden Instability}
\label{sec:day07-finding}

\textbf{\texttt{top\_region\_overlap\_score}'s 0.667 mean looks reasonably stable. It is hiding how
unstable the safety-relevant box actually is.} The metric is computed on Day 4's area-ranked top-1
region --- the same \texttt{standardize\_regions} ranking introduced in Chapter 4, which orders
candidate boxes by pixel area with no access to a true importance signal. For rules 1--3, the
area-ranked top-1 region is almost always the \textbf{worker} box, simply because a worker's
bounding box is much larger in pixels than the hard hat, harness, or guardrail the rule is actually
asking about. \texttt{top\_region\_overlap\_score} at 0.667, in other words, is substantially a
measurement of how stable Florence-2's worker-detection box is across stochastic reruns ---
informative in its own right, but not a direct answer to the question this chapter actually needs
answered: is the explanation grounded in the rule's safety object stable?

To answer that question directly, \texttt{object\_region\_overlap\_score} was added as a
\emph{separate} metric, computed specifically on \texttt{object\_boxes[0]} --- the rule's actual
safety-object box, not whichever box happens to rank first by area. The gap this exposes is large.
Rule\_1's hard-hat box overlaps only \textbf{0.499} across reruns, against rule\_4's excavator box
at \textbf{0.922} --- nearly \textbf{2$\times$} more stable. Averaged across all 163 samples,
\texttt{object\_region\_overlap\_score} comes in at \textbf{0.602}, noticeably below
\texttt{top\_region\_overlap\_score}'s 0.667 --- a smaller gap in the aggregate than the by-rule
breakdown shows, because rule\_4's excavator pulls the object-score average up even as rules 1--3
pull it down. The aggregate alone understates how unevenly that instability is distributed; the
by-rule table is where the real picture is.

\subsection{Ruling out box reordering}

An easy alternative explanation needed to be checked and ruled out before attributing this gap to
anything more interesting: could \texttt{object\_region\_overlap\_score}'s instability simply be
multiple detected boxes getting \emph{reordered} across stochastic runs, so that
\texttt{object\_boxes[0]} on one rerun picks out a genuinely \emph{different} detected box than
\texttt{object\_boxes[0]} on another rerun --- a list-indexing artifact, not evidence that any one
box is actually \emph{moving}? This was checked directly rather than assumed away: only
\textbf{4--8\%} of samples have more than one object box detected at all (rule\_1 6.3\%, rule\_2
3.8\%, rule\_3 8.2\%, rule\_4 4.0\%). With reordering only possible in a single-digit-percentage
slice of samples per rule, it cannot explain a systematic, roughly 2$\times$ gap in mean IoU that
holds across the full 24--63 samples assigned to each rule. The instability is real box movement,
not list-index noise.

\subsection{The real driver: object size and shape}

With reordering ruled out, the explanation that survives is a combination of \textbf{object size
and shape} --- the same kind of confound Chapter 6 identified in the visual-sparsity metric, here
showing up in a different metric for a related reason. Median detected box area, taken from the
sampled reruns, lines up closely with \texttt{object\_region\_overlap\_score}:

\begin{center}
\begin{tabular}{llrr}
\toprule
\textbf{Rule} & \textbf{Object} & \textbf{Median box area (px\textsuperscript{2})} & \textbf{object\_region\_overlap\_score} \\
\midrule
rule\_1 & hard hat & 570 & 0.499 \\
rule\_2 & harness & 7{,}169 & 0.533 \\
rule\_3 & guardrail & 221{,}163 & 0.615 \\
rule\_4 & excavator & 215{,}069 & \textbf{0.922} \\
\bottomrule
\end{tabular}
\end{center}

The two smallest objects, hard hat and harness, are also the two least stable; the two largest,
guardrail and excavator, are the two most stable. Hard hats and harnesses are small \emph{and}
tightly bound to a worker's body, so a few pixels of box-edge jitter under stochastic decoding
represents a large fraction of the box's own size, driving IoU down sharply even when the model is
plainly still attending to roughly the same spot on the image.

Size alone, however, does not fully explain the pattern, and the data makes that point cleanly:
\textbf{rule\_3's guardrail and rule\_4's excavator have almost identical median box area} ---
221,163~px\textsuperscript{2} versus 215,069~px\textsuperscript{2}, a difference of under 3\% ---
\textbf{yet very different stability}, 0.615 against 0.922. If size were the whole story, these two
rules should land at nearly the same \texttt{object\_region\_overlap\_score}. They do not, and the
remaining gap is best read as a \textbf{shape and rigidity} effect layered on top of size. A
guardrail is long, thin, and sometimes visually segmented (formwork panels, intermittent rails),
so small differences in where the model decides the structure's run starts or stops shift the
box's extent noticeably between reruns, even when the area enclosed stays roughly similar. An
excavator is one compact, rigid, visually distinctive object whose silhouette does not depend much
on exactly where stochastic decoding happens to place the boundary --- there is comparatively
little ambiguity about where a single excavator begins and ends, which a long, segmented guardrail
does not share. Size sets the floor on how much a few pixels of jitter can hurt IoU; shape decides
how much jitter there actually is.

\section{Worked Examples}

Figure~\ref{fig:day07-worked} shows two saved rerun overlays that make this chapter's size-and-shape
finding visually concrete, one from each end of the \texttt{object\_region\_overlap\_score} range.

\begin{figure}[h]
\centering
\begin{subcaptionblock}{0.48\linewidth}
\centering
\includegraphics[width=\linewidth]{stability/0000023_rule_1_reruns.png}
\caption{Image 23, rule~1 (hard hat): the overlaid top-1 (worker) boxes from all 3 reruns sit
almost exactly on top of each other; the hard-hat box itself, not drawn here because it loses the
area ranking to the worker box, is what actually moves
(\texttt{object\_region\_overlap\_score}~$=0.069$ for this sample).}
\end{subcaptionblock}
\hfill
\begin{subcaptionblock}{0.48\linewidth}
\centering
\includegraphics[width=\linewidth]{stability/0000079_rule_3_reruns.png}
\caption{Image 79, rule~3 (guardrail): the guardrail/formwork box is nearly identical across all 3
reruns (\texttt{object\_region\_overlap\_score}~$=0.955$).}
\end{subcaptionblock}
\caption{Rerun overlays anchoring the hidden-instability finding: a stable worker detection masking
an unstable hard-hat box (left) against a guardrail box that is stable in its own right (right).}
\label{fig:day07-worked}
\end{figure}

Image 23 is the sharpest illustration of this chapter's central point. The three reruns' top-1
boxes --- all worker boxes, since the worker is the larger object and wins the area ranking ---
overlay almost perfectly, which is exactly what feeds a healthy-looking
\texttt{top\_region\_overlap\_score}. But that overlay never draws the hard-hat box, because the
hard hat loses the area ranking every time. The hard-hat box's own overlap across the same three
reruns, \texttt{object\_region\_overlap\_score}~$=0.069$ for this specific sample, is far below
even rule\_1's already-low 0.499 mean --- a near-complete miss for the object the rule is actually
about, sitting directly underneath a top-1 score that looks fine. Image 79's guardrail box, by
contrast, needs no such unmasking: it is the object actually being scored, and it sits nearly
identically across all three reruns, \texttt{object\_region\_overlap\_score}~$=0.955$, consistent
with this chapter's account of large, rigid structures as the more stable case on both counts.

\section{Standard vs. Non-Standard Choices}

\textbf{Standard}: sampling-based stability and consistency testing --- rerun the same input under
controlled stochasticity and measure rerun-to-rerun agreement --- is itself standard practice,
rooted directly in stability as one of the six metrics defined for black-box XAI evaluation
\citep{arreche2024exai}. Switching to \texttt{do\_sample=True}, \texttt{num\_beams=1},
\texttt{temperature=0.7} specifically to avoid the vacuous deterministic-rerun case, and confirming
non-vacuousness via \texttt{answer\_agreement\_rate}~$=0.775 \neq 1.0$ before trusting the
aggregate, follows that established methodology faithfully rather than improvising around it.

\textbf{Non-standard}: introducing \texttt{object\_region\_overlap\_score} as a metric
\emph{separate from} the naive top-1-region IoU is this chapter's contribution. The six-metric
framework's stability definition does not by itself specify which region a region-overlap score
should be computed on when more than one candidate box is available per sample; the natural default
--- the top-ranked region from whatever heuristic produced the ranking --- is exactly the choice
that, here, hid a confound rather than revealing it. Computing a second, target-specific overlap
score alongside the naive one, purpose-built to unmask the confound the naive metric was shown to
hide, is a methodological refinement rather than a deviation from the framework, and is presented
here as worth keeping in the framework going forward, not as a one-off patch for this pilot alone.

\section{Implications for the Paper}
\label{sec:day07-implications}

\textbf{The central methodological recommendation from this chapter is general, not specific to
stability or to this pilot}: whenever a region-overlap or region-based stability metric is computed
on top of a ranking heuristic rather than on a metric's own true target object, the paper should
report \emph{both} numbers --- the ranking heuristic's pick and the metric's own target object ---
rather than the ranking heuristic's pick alone. A single aggregate computed only on ``whichever
region a heuristic ranks first'' can look acceptable while masking a target object that is
substantially less stable, exactly as \texttt{top\_region\_overlap\_score}'s 0.667 masked
\texttt{object\_region\_overlap\_score}'s 0.602 mean and, more sharply, rule\_1's 0.499. This ties
directly to the area-ranking limitation already flagged in
Chapter~\ref{chap:day04}\footnote{Chapter 4's area-ranking heuristic, \texttt{standardize\_regions},
ranks candidate boxes by pixel area as a stand-in for true importance, and that chapter's own
implications section already forward-references this gap by name.}: any metric built on top of
that ranking inherits its blind spot, and this chapter is the second of (at least) two places in
the pilot where that inheritance produces a visible, quantified effect rather than a theoretical
concern.

\textbf{The same underlying pattern recurs here that Chapter 6 identified for visual sparsity, and
the recurrence should be stated plainly rather than treated as coincidence}: metrics computed on
small, body-worn PPE items (hard hat, harness) come out systematically worse than the same metrics
computed on large, rigid objects (excavator), for reasons that are at least partly mechanical ---
box size and shape interacting with IoU and with a 24$\times$24-patch attention grid --- rather than
purely reflecting how good the underlying explanation actually is. Chapter 6 showed this for
attention concentration; this chapter shows it for rerun stability; \textbf{Chapter 11's
cross-metric analysis is the natural place to bring the two findings together} rather than letting
rule\_1 and rule\_2's weaker numbers be read in isolation, in either chapter, as evidence of worse
explanations specifically for hard-hat and harness judgments. The irony is worth stating outright in
the paper's discussion: hard hat and harness are exactly the PPE items for which a \emph{stable}
explanation matters most for real-world trust in a safety system, and they are the least stable by
this measure, for reasons that trace back to object geometry rather than to any deficiency in how
Florence-2 reasons about them.
